% =========================================================================
% Section V -- Experiments.
% Numerical sources:
%   C:/Users/rtes/Desktop/AMR-DFJSP/experiments/full_benchmark
%   C:/Users/rtes/Desktop/AMR-DFJSP/experiments/surrogate_fidelity
% =========================================================================

\section{Experiments}
\label{sec:experiments}

\subsection{Protocol and Comparison Methods}

We keep the $20\times19$ layout, the homogeneous $m=16$ fleet, the station
configuration, and the fixed executor $\Phi$ unchanged, and vary workload
through $n\in\{20,40,60,80,100\}$. Job class, inbound station, and outbound
station are sampled uniformly, and all jobs are released at zero. Training uses
5,000 independently generated 60-job instances. Fixed validation and test sets
are mutually disjoint and are also disjoint from training; the test set contains
100 instances at each workload. Before test evaluation, every learned policy is
selected by the same executor-based validation score,
$C_{\max}^{\Phi}+1000\nu_\Phi$, so failed decodes cannot be preferred to
completed schedules.

The dispatching-rule library crosses ten job priorities with six AMR rules,
giving 60 deterministic combinations. The job rules cover FIFO, short- and
long-processing-time priorities, distance and station workload, predicted
completion, material matching, milk-run batching, and random choice; the AMR
rules cover availability, distance, predicted completion, load, material
matching, and random choice. For the baseline comparison, one combination is
selected on the validation split separately for each $n$ and frozen before test
evaluation. The genetic algorithm (GA) uses a population of 200 for 150
generations followed by collision-aware local improvement.

The learned-policy experiment crosses two observations with three terminal
returns. The core observation $O_{\mathrm{core}}$ removes the projected
congestion channels of Section~\ref{sec:arch-cost}, whereas
$O_{\mathrm{cong}}$ retains them. Completed rollouts are priced by the raw
idealized makespan $\widetilde C$, the executor makespan $C_{\max}^{\Phi}$, or
the calibrated surrogate makespan $\hat\Psi$. Configuration names concatenate
the return and observation; Surr--Cong is the proposed configuration. All six
arms use the same architecture, data, and 4,000-epoch PPO budget.

All arms use the same executor-based model-selection rule, preventing
evaluator-specific checkpoint selection.

At inference we report one deterministic greedy rollout and best-of-eight
sampling, which selects the routable candidate with the lowest executed
makespan. Primary comparisons average three training seeds within each
instance. The experimental unit is the instance, method comparisons are paired,
and 95\% intervals are $1.96s/\sqrt{N}$. Failed decodes receive the finite
score above during training and model selection and are reported as failures.
Because that score is not an elapsed completion time, makespan averages include
completed schedules only.

\subsection{Evaluator Fidelity}
\label{sec:surrogate-results}

We first test whether $\widetilde\Phi$ preserves schedule order rather than only
absolute value. This broad comparison uses 19 candidates per instance spanning
rules, GA, and learned policies; more than 99.6\% are executor-routable.
Table~\ref{tab:ranking} excludes failed schedules and uses the metrics of
Section~\ref{sec:evaluators}. Here $n$ is the number of jobs and $\tau_b$ is the
per-instance, tie-corrected Kendall coefficient with respect to $\Phi$.

\begin{table}[!ht]
\caption{Idealized--executor ranking.}
\label{tab:ranking}
\centering
\scriptsize
\setlength{\tabcolsep}{3.0pt}
\begin{tabular}{@{}rrrr@{}}
\toprule
$n$ & $\tau_b$ & Inv. (\%) & Reg. (\%) \\
\midrule
20  & $0.26\pm0.04$ & $36.25\pm1.93$ & $16.08\pm2.73$ \\
40  & $0.13\pm0.03$ & $43.22\pm1.54$ & $38.62\pm3.43$ \\
60  & $0.23\pm0.03$ & $38.40\pm1.27$ & $30.63\pm3.79$ \\
80  & $0.30\pm0.03$ & $34.96\pm1.30$ & $27.51\pm7.20$ \\
100 & $\mathbf{0.37\pm0.02}$ & $\mathbf{31.39\pm1.20}$ & $\mathbf{9.41\pm4.04}$ \\
\bottomrule
\end{tabular}
\end{table}

The abstraction misranks schedules at every workload. At the 60-job training
size, $38.4\%$ of comparable pairs are inverted and selecting its minimum
incurs $30.6\%$ executed regret. Although ranking partially recovers at 100
jobs, $\tau_b=0.37$ and $9.4\%$ regret show that value calibration alone does
not guarantee decision fidelity.

We next compare $\widetilde C$ and $\hat\Psi$ with $\Phi$ on the same
19-candidate field: 12 competitive rules, all six factorial policy cells, and
GA. The rules combine FIFO, SPT, LPT, milk-run, material-matching, or
earliest-completion job priority with the earliest-available or
earliest-completion AMR rule. They occupy mean ranks 3--30 in the full 60-rule
benchmark; excluding clearly dominated random, nearest-station, and
least-congested rules prevents mechanical inflation of rank agreement. The
policies and GA also prevent $\hat\Psi$ from being tested only on schedules
constructed by rules that already use it. Table~\ref{tab:surrogate} reports
within-instance decision error and the total time to score all 19 schedules of
one instance; candidate generation is excluded.

\begin{table*}[!t]
\caption{Evaluator fidelity and scoring time.}
\label{tab:surrogate}
\label{tab:evaluator-time}
\centering
\scriptsize
\setlength{\tabcolsep}{3.0pt}
\begin{tabular}{@{}r|c|cccc|cccc@{}}
\toprule
& Executor $\Phi$ & \multicolumn{4}{c}{Idealized $\widetilde C$}
& \multicolumn{4}{c}{Surrogate $\hat\Psi$} \\
\cmidrule(lr){2-2}\cmidrule(lr){3-6}\cmidrule(lr){7-10}
$n$ & $t$ (ms) & $\tau_b$ & Inv. (\%) & Reg. (\%) & $t$ (ms)
    & $\tau_b$ & Inv. (\%) & Reg. (\%) & $t$ (ms) \\
\midrule
20  & 664.66 & $0.26\pm0.04$ & 36.25 & 16.08
    & \textbf{1.18} & $\mathbf{0.69\pm0.03}$
    & \textbf{14.91} & \textbf{3.62} & 2.79 \\
40  & 1279.83 & $0.13\pm0.03$ & 43.22 & 38.62
    & \textbf{2.03} & $\mathbf{0.73\pm0.02}$
    & \textbf{13.04} & \textbf{2.69} & 5.17 \\
60  & 1915.73 & $0.23\pm0.03$ & 38.40 & 30.63
    & \textbf{2.84} & $\mathbf{0.75\pm0.02}$
    & \textbf{12.28} & \textbf{3.19} & 7.58 \\
80  & 2636.45 & $0.30\pm0.03$ & 34.96 & 27.51
    & \textbf{3.65} & $\mathbf{0.77\pm0.02}$
    & \textbf{11.06} & \textbf{2.65} & 10.01 \\
100 & 3243.23 & $0.37\pm0.02$ & 31.39 & 9.41
    & \textbf{4.31} & $\mathbf{0.81\pm0.01}$
    & \textbf{9.50} & \textbf{1.28} & 12.22 \\
\bottomrule
\end{tabular}
\end{table*}

Decision metrics are computed within each instance before aggregation, so the
effective sample size is 100 rather than 1,900 schedules. Kendall's $\tau_b$
and Inv. handle ties as defined in Section~\ref{sec:evaluators}; selection
regret uses the candidate order fixed before test evaluation. The executor
rejects 26 of the 9,500 schedules (0.27\%); decision metrics use the common
executor-routable candidate set, whereas each time observation totals all 19
evaluator calls for one instance before the mean is computed.

\subsection{Terminal Return and Congestion Observability}

Table~\ref{tab:factorial} gives the single-seed best-of-eight factorial, with
surrogate configurations in the rightmost block. Both executor and surrogate
returns substantially improve on the idealized-return controls; the comparison
of greatest practical interest is therefore Surr--Cong versus Exec--Cong.

\begin{table}[!ht]
\caption{Best-of-eight factorial results.}
\label{tab:factorial}
\centering
\scriptsize
\setlength{\tabcolsep}{1.8pt}
\resizebox{\columnwidth}{!}{%
\begin{tabular}{@{}rcc|cc|cc@{}}
\toprule
& \multicolumn{2}{c}{Ideal return} & \multicolumn{2}{c}{Executor return}
& \multicolumn{2}{c}{Surrogate return} \\
\cmidrule(lr){2-3}\cmidrule(lr){4-5}\cmidrule(lr){6-7}
$n$ & Core & Cong & Core & Cong & Core & Cong \\
\midrule
20  & $192.70 \pm 3.60$ & $222.72 \pm 4.20$
    & $160.18 \pm 2.96$ & $161.14 \pm 2.93$
    & $\mathbf{160.16 \pm 2.65}$ & $162.00 \pm 2.92$ \\
40  & $378.69 \pm 6.05$ & $325.48 \pm 4.39$
    & $246.58 \pm 2.57$ & $242.15 \pm 3.08$
    & $253.51 \pm 3.07$ & $\mathbf{239.23 \pm 3.25}$ \\
60  & $603.03 \pm 7.03$ & $419.62 \pm 4.74$
    & $338.48 \pm 2.11$ & $326.07 \pm 3.05$
    & $351.88 \pm 2.52$ & $\mathbf{319.89 \pm 3.09}$ \\
80  & $799.95 \pm 6.34$ & $524.35 \pm 4.46$
    & $436.05 \pm 2.61$ & $416.56 \pm 3.16$
    & $450.99 \pm 3.21$ & $\mathbf{403.43 \pm 3.83}$ \\
100 & $968.25 \pm 8.27$ & $630.15 \pm 5.51$
    & $534.46 \pm 2.49$ & $504.30 \pm 3.47$
    & $550.14 \pm 3.41$ & $\mathbf{488.50 \pm 3.73}$ \\
\bottomrule
\end{tabular}%
}
\end{table}

Under $O_{\mathrm{cong}}$, Surr--Cong is better than Exec--Cong from 40
through 100 jobs, whereas the same pattern is not observed under
$O_{\mathrm{core}}$. Best-of-eight strengthens the result: Surr--Cong is
better at all five workloads by $2.0$--$15.5$ ticks. Thus, a higher-fidelity
executor return does not automatically yield the better policy under
surrogate-based state transitions.

\subsection{Comparison with Rules and Genetic Search}

Table~\ref{tab:baseline-performance} compares Surr--Cong with baselines fixed
without using the test set and reports inference time. Greedy performance is
mixed, whereas best-of-eight reduces rule makespan by 4.1--9.1\% at every
workload. It is 3.1\% worse than GA at $n=20$, but 3.6--7.9\% better from
$n=40$ through 100.

\begin{table}[!ht]
\caption{Executed makespan and scheduling time.}
\label{tab:baseline-performance}
\label{tab:inference-time}
\centering
\scriptsize
\setlength{\tabcolsep}{2.0pt}
\resizebox{\columnwidth}{!}{%
\begin{tabular}{@{}rcc|cc|cc|cc@{}}
\toprule
& \multicolumn{2}{c}{Rule} & \multicolumn{2}{c}{GA}
& \multicolumn{2}{c}{Surr-G} & \multicolumn{2}{c}{Surr-8} \\
\cmidrule(lr){2-3}\cmidrule(lr){4-5}\cmidrule(lr){6-7}
\cmidrule(lr){8-9}
$n$ & $M$ & $t$ & $M$ & $t$ & $M$ & $t$ & $M$ & $t$ \\
\midrule
20  & $166.22 \pm 2.82$ & \textbf{0.03}
    & $\mathbf{153.95 \pm 2.84}$ & 15.73
    & $168.65 \pm 3.16$ & 0.10
    & $158.78 \pm 2.93$ & 1.06 \\
40  & $258.38 \pm 3.78$ & \textbf{0.10}
    & $246.34 \pm 3.00$ & 37.21
    & $249.39 \pm 3.83$ & 0.35
    & $\mathbf{237.54 \pm 3.14}$ & 3.17 \\
60  & $349.61 \pm 3.57$ & \textbf{0.23}
    & $338.68 \pm 2.98$ & 65.85
    & $334.75 \pm 3.34$ & 0.77
    & $\mathbf{317.95 \pm 2.65}$ & 6.63 \\
80  & $426.24 \pm 7.22$ & \textbf{0.57}
    & $435.85 \pm 4.02$ & 101.86
    & $422.39 \pm 3.94$ & 1.36
    & $\mathbf{401.38 \pm 3.40}$ & 11.83 \\
100 & $505.98 \pm 8.02$ & \textbf{0.89}
    & $525.85 \pm 4.07$ & 144.93
    & $513.23 \pm 4.46$ & 2.15
    & $\mathbf{485.11 \pm 3.11}$ & 18.70 \\
\bottomrule
\end{tabular}%
}
\end{table}

Best-of-eight produces no routing failures in the reported Surr--Cong results;
greedy decoding fails in 7 of 1,500 seed--instance evaluations. As specified in
the protocol, these failures are not treated as completed makespans.

\subsection{Computation Time}

We separate evaluator cost from end-to-end scheduling cost.
Table~\ref{tab:evaluator-time} reports the mean elapsed time
per test instance for scoring its 19 candidates; it excludes candidate
generation.

The scheduling-time columns of Table~\ref{tab:inference-time} report serial,
single-process results on an NVIDIA GeForce RTX 5080. The common final
reporting call to $\Phi$ is excluded, while executor calls used for
best-of-$K$ selection are included.

Training required approximately 55--74 hours per run under shared GPU load.
Because contention differed and validation still invokes $\Phi$, these values
document resource use but do not support a controlled whole-training speedup
claim.
